\documentclass[main.tex]{subfiles}


\begin{document}

%from pagenumber to up
% \phantomsection 
% \hypertarget{toc}{}

 

 

 
   

\section{ЭДС индукции в движущемся проводе}


Пусть прямой провод перемещается в постоянном магнитном поле (рис.\,\ref{fig:p161}).

    \begin{figure}[H]
    \centering
\includegraphics[]{el/43mcond_pic1.pdf}%
\caption{Движение провода в магнитном поле}
\label{fig:p161}
\end{figure}

Провод~$CD$ движется в магнитном поле~$B$ в плоскости, перпендикулярной линиям поля, поступательно так, что его скорость~$\vec v$ образует прямой угол с осью провода ($\vec v \perp CD$).
Опыт показывает, что в этом случае между концами провода возникает \emph{напряжение.} 
% которую можно измерить \emph{неподвижным} вольтметром, подключенным к этим концам.
В проводе как бы появляется <<невидимая батарейка>> с ЭДС индукции (без внутреннего сопротивления).

\textbf{ЭДС индукции в движущемся проводе} при его движении в магнитном поле равна:
\begin{equation}
    \label{43mcond_e1}
    \mathscr E_i=Bvl\cos\alpha,
\end{equation}
где $B$~--- индукция магнитного поля, $v$~--- скорость провода, $l$~--- длина провода (находящегося в магнитном поле; $\vec v\perp l$), $\alpha$~--- угол между вектором~$\vec B$ и перпендикуляром к плоскости движения провода.

На рис.\,\ref{fig:p162} изображен пример к определению ЭДС индукции в движущемся проводе.

    \begin{figure}[H]
    \centering
\includegraphics[]{el/43mcond_pic2.pdf}%
\caption{К определению ЭДС индукции в движущемся проводе}
\label{fig:p162}
\end{figure}

В магнитном поле с индукцией~$\vec B$ поступательно движется провод длиной~$l$ со скоростью~$\vec v$ (причем $\vec v\perp l$), при этом вектор~$\vec B$ образует угол~$\alpha$ с перпендикуляром~П к плоскости движения. Формула~\eqref{43mcond_e1} дает ЭДС индукции в этом проводе: $\mathscr E_i=Bvl\cos\alpha$. 






























 
\end{document}